\chapter{Topic}
\label{topic}
This project is part of the course Speech Recognition (ELEC-E5150) of the year 2026 at Aalto University.

The goal of this project is to build a Speech Recognition System for an artificial dataset containing ghibberish speech samples of Esperanto.
Esperanto is a planned language that was invented by the polish doctor Ludwik Lejzer Zamenhof at the end of the 19th century to serve as a language for simplified international communication.
To enable as many as people as possible learn it easily, it was especially designed to have only very few grammatical rules without any exceptions and relies on a strong and fixed coupling between the written and spoken langage.
In recent years the failing to stablish Esperanto as a global language became obvious with the rise and acceptance of English for international communication.
Nevertheless, there are Esperanto speakers on every continent today.

Throughout the remainings of this project we will develop an automated system to transcribe Esperanto to its textual representation.
The used dataset contains no real speech samples but ghibberish ones and is part of the Common Voice Corpus. 